\documentclass[runningheads]{llncs}
\usepackage{graphicx}
\usepackage{booktabs}
\usepackage{multirow}
\usepackage{microtype}
\usepackage{xcolor}
\usepackage[hidelinks]{hyperref}
\emergencystretch=1em
\newcommand{\method}{CEARF-N}
\begin{document}
\title{Rejection-Complete Validation Gating for Sparse Session Recommendation}
\titlerunning{Rejection-Complete Validation Gating}
\author{Anonymous Author(s)}
\authorrunning{Anonymous Author(s)}
\institute{Anonymous Institution}
\maketitle

\begin{abstract}
Sparse, large-catalogue next-item recommendation makes it unclear when neural
sequence evidence is more reliable than explicit local memory. We introduce
\emph{rejection-complete validation gating}: every optional source, neural
mechanism, fusion weight, and router is paired with an explicit null or simpler
alternative, and complexity is admitted only by held-out utility. We instantiate
the principle in \method{}, which combines transition, session-neighbour,
popularity, and compact neural rankings using weighted reciprocal-rank
fusion~\cite{cormack2009rrf}. Under full-catalogue evaluation on two sparse
Amazon domains and Diginetica, the selected fusion improves over both of its
own endpoints on all three domains. The metadata-equipped system obtains the
strongest Recall@20 among evaluated ID-only external baselines on both Amazon
domains; semantic-initialized NARM overtakes it under matched side information,
and ID-only NARM is significantly stronger on Diginetica. A fourth,
weak-metadata Tmall stress test provisionally selects CEARF-N+STAN, whereas
an expanded expert family selects CEARF-N+STAN+NARM on Video Games, excludes
NARM on Baby Products, and selects STAN+NARM without CEARF-N on Diginetica.
Thus the contribution is not a universal winner or a new fusion formula, but a
falsifiable selection method that identifies when heterogeneous evidence earns
inclusion and exposes when apparent gains come from side information or a weak
primary model.
\keywords{next-item recommendation \and sparse recommendation \and rank fusion
\and validation-based model selection \and hybrid recommender systems}
\end{abstract}

\section{Introduction}
Next-item recommendation ranks the item most likely to follow an observed
interaction sequence. The task includes anonymous sessions as well as
session-like prefixes extracted from persistent user histories. Although these
settings share a ranking objective, they do not share the same data-generating
process; we therefore evaluate them separately and avoid claims based on
averages across domains.

Sparse, large-catalogue data creates three recurring failures. First, neural
sequence models allocate parameters to items whose representations may receive
too few updates to become reliable. Recent state-space, semantic-embedding,
and long-tail-oriented systems improve efficiency or representation quality,
but they do not remove this basic long-tail constraint~\cite{sigma2025,
llm2rec2025,longtail2026}. Second, neighbourhood and transition methods remain
strong under careful evaluation, and a 2026 audit shows that local graph
evidence alone can solve a surprising fraction of widely used sequential
benchmarks~\cite{popstrat2025,graphshortcut2026}. Their weakness is not lack
of signal but lack of adaptation: the same heuristic is applied whether a
query has rich collaborative support or almost none. Third, score-level
hybrids combine transition counts, neighbourhood similarities, and neural
logits whose scales and tails are fundamentally different. A weight that works
in one domain can therefore fail in another for purely numerical reasons.

The deeper methodological problem is that recommender-system development
usually asks validation to tune a proposed complex model, not whether that
complexity should exist. Reproducibility studies repeatedly find carefully
tuned neighbourhood or linear methods competitive with recent neural
proposals~\cite{ferraridacrema2019progress,ludewig2021empirical}. Existing
hybrids can improve accuracy by combining recurrent and neighbourhood
models~\cite{jannach2017hybrid}, but a mandatory hybrid cannot return the
scientifically meaningful answer ``do not add this component.'' We address
this missing falsifiability rather than claiming that simple recommendation
itself is a new observation.

These observations motivate a different architectural prior: use explicit
memory as the stable base, express each component only through rank, and admit
neural evidence only when held-out validation supports it. This is the
principle behind \method{}, a memory-first rank-fusion system with a compact
neural residual. The fusion primitive is deliberately standard. The new object
is instead a \emph{rejection-complete} configuration space: every proposed
mechanism has a null endpoint, every router competes with a global constant,
and exact router-selection ties resolve toward the simpler system.
Consequently, the procedure can falsify a design idea rather than merely tune
its strength.

The need for such a reference point has increased. Popularity-stratified
analysis shows that neighbourhood and neural recommenders win on different
regions of the same data~\cite{popstrat2025}, while algorithm-selection work
shows that per-instance routing closes only part of the gap to an
oracle~\cite{algselect2025}. At the same time, learned ensembles and LLM
routing systems can carry substantial serving and energy cost~\cite{
flame2026,dynamicllm2026,energy2026,green2025}, and evaluation defaults can
change apparent winners~\cite{defaults2025,seqsplit2026}. A lightweight hybrid
is therefore useful not only as a practical system, but as a cost floor that
more elaborate approaches should exceed.

This paper asks three questions: (RQ1) which optional mechanisms survive when
validation may select their null states; (RQ2) how much of the Amazon ordering
is caused by semantic metadata rather than fusion; and (RQ3) whether rank
fusion adds value beyond combining strong external experts. We address the
third question with a NARM-containing control family, reported as a
single-artifact diagnostic pending matched-seed replication.

This paper makes four contributions:

\begin{enumerate}
\item \textbf{Rejection-complete validation gating.} We formulate selection
      so that every optional component has an explicit off state and every
      adaptive router competes with a constant alternative. Held-out utility
      therefore decides both inclusion and exclusion; router complexity is
      used only to break exact ties.
\item \textbf{A rank-space instantiation.} Three explicit memories and a
      compact neural residual are combined without cross-component score
      calibration. The family includes memory-only and neural-only boundaries,
      making ``fusion helps'' directly falsifiable.
\item \textbf{A controlled full-catalogue study.} Three matched-seed domains,
      eight external baselines, paired inference, constituent ablations, a
      no-metadata study, and a semantic-fairness rerun separate the system-level
      result from the fusion-level result. A fourth Tmall diagnostic tests
      transfer without promoting provisional evidence to the main comparison.
\item \textbf{Metadata-confound accounting.} Removing semantic initialization
      and then granting the same initialization to NARM reverses the Amazon
      ordering. We report this reversal as a primary finding rather than
      attributing it to the fusion architecture.
\end{enumerate}

Our claims are deliberately bounded. We report the best result within one
unified protocol, not a literature-wide state of the art. The strongest
evidence is that \method{} combines its own constituents successfully across
all three domains; the Amazon advantage over external baselines additionally
depends on rich metadata; and broader expert fusion is helpful in some regimes
but not all.

\section{Related Work}
\subsection{Sparse sequential recommendation}
Modern sequential recommenders increasingly use state-space, semantic, and
tail-aware architectures to improve long-context efficiency or sparse-item
coverage~\cite{sigma2025,llm2rec2025,longtail2026}. Long-tail accuracy
nevertheless remains unresolved: stronger sequence modeling does not eliminate
the core problem that rare items receive weak or noisy evidence. \method{}
acts earlier in the pipeline. Instead of requiring every tail item to obtain a
strong learned representation, it allows explicit transition and session
memories to supply evidence directly.

Recent evaluation work also changes how progress should be interpreted.
Popularity-stratified analysis finds carefully tuned local methods competitive
with neural systems on some strata~\cite{popstrat2025}. More sharply, a 2026
audit demonstrates that a simple graph heuristic explains much of the
performance on several sequential benchmarks~\cite{graphshortcut2026}. This
does not make neural modeling unnecessary; it means a neural model should be
required to add value beyond local memory rather than receive that assumption
by default.

\subsection{Hybrid ranking and model selection}
Complementarity between recommenders is often hidden by aggregate metrics.
Popularity-stratified modeling shows that session-neighbourhood and neural
methods dominate different item strata~\cite{popstrat2025}. More elaborate
combinations include compact ensemble distillation and dynamic LLM
routing~\cite{flame2026,dynamicllm2026}. \method{} takes the opposite route:
it retains heterogeneous evidence sources but prevents their scores from
interacting directly. Only ranks are fused, and every profile is selected from
a finite validation grid.

Per-query selection is attractive but difficult. A 2025 meta-learning study
reports improvement over a single global algorithm while closing only a small
part of the oracle gap~\cite{algselect2025}. We therefore treat routing as
secondary rather than foundational: a global or regime-level $\beta$ remains a
valid outcome, and finer routing survives only if it wins on a disjoint
validation fold.

\subsection{Reliability, side information, and cost}
Semantic item representations are increasingly used to address sparse
interaction data~\cite{llm2rec2025}. Because semantic information can make an
unfair comparison appear architectural, we report both a no-metadata ablation
and a fairness rerun that gives external baselines the same initialization.
Reliable comparison also requires more than using the same metric. Evaluation
defaults can change apparent winners~\cite{defaults2025,seqsplit2026}, and
energy-aware studies show that ensemble gains can carry disproportionate
computational cost~\cite{energy2026,green2025}. \method{} is designed as an
auditable reference system whose selected configuration can be reproduced on a
laptop.

\section{Method}
A query is the most recent interaction context, truncated to eight items.
Offline, \method{} builds three memories, trains one compact residual model,
and selects all discrete choices on validation data. Online, each component
emits a list of at most 120 candidates; no raw score crosses component
boundaries.

\subsection{Evidence memories}
\textbf{Transition memory.} A directed graph counts item transitions within a
four-step window. Each source retains its 200 most frequent targets. For
context item $s$ at recency age $a$ and candidate $t$,
\begin{equation}
T(t)=\sum_a e^{-0.35a}\frac{\log(1+c(s_a,t))}{\sqrt{1+\mathrm{freq}(t)}}.
\end{equation}
The decay privileges recent evidence, while the frequency denominator prevents
globally popular targets from dominating every context.

\textbf{Session memory.} Candidate training sessions are retrieved by
IDF-weighted, recency-decayed overlap:
\begin{equation}
O(q,j)=\sum_a \mathrm{idf}(i_a)\cdot e^{-0.20a}\cdot \mathbf{1}[i_a\in j].
\end{equation}
The top 120 sessions vote for candidates. Items appearing after the last match
receive weight 1.0; earlier items receive 0.25. Votes are discounted by
square-root distance and normalized by unique session length.

\textbf{Popularity memory.} Global popularity provides a weak third ranked list
inside the validation-selected memory profile and also pads a fused list
containing fewer than twenty candidates. It is therefore both low-weight
evidence and conservative fallback, but it never dominates context-dependent
signals because its validation-selected profile weight remains small.

\subsection{Rank-space fusion}
For memory list $i$, profile weight $w_i$, and reciprocal-rank constant
$k=20$, the memory score is
\begin{equation}
M(t)=\sum_i \frac{w_i}{20+\mathrm{rank}_i(t)}.
\end{equation}
Rank space makes the components comparable without calibrating transition
counts, overlap scores, or frequencies. Six fixed profiles over transition,
session, and popularity evidence are considered on validation.

A multiplicative agreement bonus was evaluated but changed Recall@20 by
.00000 on every domain. Because weighted reciprocal-rank fusion already
rewards items supported by multiple lists, the extra term is retained only to
match the locked implementation and is interpreted as inert.

\subsection{Neural residual}
PASGR is a single-layer GRU with hidden size 64 and dropout .2. Item
embeddings are initialized from a 128-dimensional semantic matrix. On Amazon,
the matrix is obtained by truncated singular-value decomposition of TF--IDF
vectors built from titles, categories, features, descriptions, and brands. On
Diginetica, product names are mostly numeric tokens, reducing the usefulness of
semantic content and creating a natural weak-metadata contrast.

The residual is trained with cross-entropy. Graph transport, contrastive loss,
in-batch loss, and prototype transport are optional validation-gated
components. The prototype transport variant is rejected on every domain; it is
therefore an evaluated alternative, not part of the final deployed system.

Memory and neural rankings are combined by
\begin{equation}
F(t)=\frac{1-\beta}{20+\mathrm{rank}_M(t)}+\frac{\beta}{20+\mathrm{rank}_N(t)},
\end{equation}
where $\beta\in\{0,.05,\ldots,1\}$. $\beta=0$ is a valid outcome and yields a
fully training-free prediction path.

\subsection{Validation-gated selection}
All choices maximize $.5\cdot\mathrm{Recall}@6+.5\cdot\mathrm{Recall}@20$ on
5,000 validation queries:
\begin{enumerate}
\item six memory profiles, separately for contexts of length at most two and
      longer contexts;
\item a 16-cell neural sweep over graph transport, prototype transport,
      contrastive loss, and in-batch loss;
\item the 21-value $\beta$ grid;
\item regime-level, bucketed, and continuous $\beta$ routers.
\end{enumerate}

\paragraph{Definition (relative rejection-completeness).}
For validation queries $V$, a declared candidate family $\mathcal{C}_j$ at
stage $j$ is rejection-complete relative to an optional mechanism $m$ if it
contains a null candidate $c_j^0$ that is identical at that stage except that
$m$ is absent. A stagewise protocol is rejection-complete relative to its
declared families when each optional mechanism has such a null candidate.
Selection at stage $j$ is
\begin{equation}
c_j^\star=\arg\max_{c\in\mathcal{C}_j} U_V(c).
\end{equation}
The qualifier ``relative'' is essential: the protocol makes no claim about
components outside the families declared before evaluation. In this
instantiation, the families contain $\beta=0$,
$\beta=1$, an off state for every neural switch, and a constant alternative to
each adaptive router. This differs from selecting weights inside a mandatory
hybrid: the validation outcome may reject the neural path, the memory path, an
auxiliary loss, or adaptation itself. We report those rejections as results.

\paragraph{Proposition.}
At every stage, $U_V(c_j^\star)\geq U_V(c_j^0)$ for every declared null
candidate $c_j^0\in\mathcal{C}_j$.
\emph{Proof.} The inequality follows directly from the argmax definition.
Consequently, a test loss against the corresponding null candidate is an
observable validation-to-test generalization failure, not a permitted
post-hoc design choice.

Router-family selection is nested within validation. Eighty percent of
validation queries fit and calibrate each router family; the disjoint twenty
percent selects the family, with simpler routers winning exact ties. The
selected family is then refit on all validation queries. Test labels are never
consulted during selection.

For Diginetica, each validation target and its incoming transition are removed
before any validation memory or model is constructed. Full training sessions
are restored only after selection for the single test fit. This prevents a
held-out next item from leaking through the transition graph.

The recency decays (.35 and .20), graph and candidate caps (200 and 120),
vote weights (1.0 and .25), and rank constant ($k=20$) are fixed
a priori and lie outside the declared candidate families. Rejection-completeness
is therefore relative to the switches and alternatives listed above, not to
every numerical constant in the implementation.

\section{Experimental Setup}
\subsection{Datasets and protocol}
\begin{table}[t]
\caption{Unified full-catalogue protocol.}
\label{tab:data}\centering\small
\begin{tabular}{lrrrr}
\toprule Dataset & Items & Training sequences & Test queries & Short contexts \\
\midrule
Amazon Video Games & 25,613 & 94,762 & 94,762 & 0\% \\
Amazon Baby Products & 36,014 & 150,777 & 150,777 & 0\% \\
Diginetica & 43,098 & 186,670 & 60,858 & 45\% \\
Tmall$^\dagger$ & 40,728 & 351,268 & 25,898 & -- \\
\bottomrule
\end{tabular}
\par\raggedright\scriptsize $^\dagger$Exploratory single-artifact transfer
diagnostic; excluded from matched-seed significance claims.
\end{table}

The Amazon domains are derived from the public Amazon review corpus~\cite{
mcauley2015amazon} and use leave-last-out next-item prediction with
exclude-seen ranking. Diginetica uses the CIKM Cup split and allows repeat
consumption. Tmall follows the public COTREC preprocessing of the IJCAI-15
shopping-log data~\cite{xia2021cotrec}.
Evaluation is full-catalogue: every candidate item is ranked and no sampled
negatives are introduced. Recall@20 is the primary measure because the target
objective is coverage within a short recommendation list. All stochastic
systems in the main comparison use matched seeds 42, 123, and 456. Tmall uses
351,268 anonymous training sessions and a repeat-consumption protocol; it has
no usable text metadata and is reported only as a seed-42 transfer diagnostic.

The benchmark is intentionally small but heterogeneous. Video Games and Baby
Products contain 625,062 and 939,529 training interactions, mean sequence
lengths 6.60 and 6.23, and about 27.7\% of items occurring at most five
times. Diginetica contributes 906,140 training interactions, shorter
session-native contexts (mean length 4.85), and a lower extreme-tail share of
19.3\%. We therefore interpret the suite as three distinct operating regimes
rather than as repeated measurements of one regime.

\subsection{Baselines and tuning}
\textbf{Baselines.} The external comparison contains five neural
baselines---GRU4Rec~\cite{hidasi2016gru4rec},
SASRec~\cite{kang2018sasrec}, NARM~\cite{li2017narm},
SR-GNN~\cite{wu2019srgnn}, and SIGMA~\cite{sigma2025}---and three non-neural
baselines---V-SKNN~\cite{ludewig2018session}, STAN~\cite{garg2019stan}, and
a transition-graph retriever. Neural baselines train for at most twelve epochs
over fifteen validation-selected configurations per seed. The reported
checkpoint is the epoch maximizing the same validation utility used by
\method{}. V-SKNN, STAN, and transition retrieval are deterministic. STAN uses
deterministic sequence order in place of wall-clock recency because the unified
loader does not retain timestamps.

The twelve-epoch neural budget is a reproducible short-paper constraint rather
than a claim of exhaustive baseline tuning. In particular, under-converged
SASRec results should be read as an internal reference point, not as a
literature-wide estimate of SASRec's best possible accuracy.

\method{} evaluates roughly seventy finite choices, but most profile, $\beta$,
and router comparisons reuse fixed candidate rankings. Only the sixteen-cell
neural-component sweep retrains PASGR. Reproducing a selected \method{}
configuration takes approximately six minutes; the one-off neural search costs
approximately eighty minutes per domain and seed.

\subsection{Statistical inference}
Pairwise differences use fingerprint-matched queries. A query-clustered paired
bootstrap with 20,000 repetitions resamples queries while preserving the three
matched predictions associated with each query; aggregate randomization uses
query-level sign flips. Exact per-seed McNemar tests provide a secondary
check. Diginetica source-session IDs are unavailable in the released HID
artifact, so a coarser session-level bootstrap cannot be recovered and is
disclosed as a limitation.

All experiments run on an Apple M2 Pro with 12 CPU cores and 32 GB memory.
The memory path is single-process CPU code; PASGR uses the available Metal
backend during training.

\section{Results}
\subsection{Main comparison}
\begin{table}[t]
\caption{Recall@20. Stochastic methods are shown as mean (std); deterministic
baselines have zero seed variance. Bold marks a clear best point estimate; no
Baby Products entry is emphasized because its leading gap is small relative
to seed variation.}
\label{tab:main}\centering\scriptsize
\setlength{\tabcolsep}{2.2pt}
\begin{tabular}{lrrrrrrrrr}
\toprule
Dataset & V-SKNN & STAN & Trans. & GRU & SAS & NARM & SR-GNN & SIGMA & \method{} \\
\midrule
Video Games & .11936 & .12115 & .04223 & .10368 (.00089) & .03794 (.00014) & .13705 (.00081) & .12267 (.00060) & .08644 (.00215) & \textbf{.14685 (.00038)} \\
Baby Products & .05038 & .04944 & .00735 & .04501 (.00056) & .03163 (.00357) & .02990 (.00018) & .05208 (.00127) & .02473 (.00059) & .05346 (.00397) \\
Diginetica & .50506 & .51502 & .40102 & .41785 (.00108) & .27097 (.00204) & \textbf{.53406 (.00052)} & .49267 (.00468) & .37219 (.01094) & .51681 (.00268) \\
\bottomrule
\end{tabular}
\end{table}

\method{} ranks first against the evaluated ID-only external baselines on both
Amazon domains. Relative to the strongest such baseline, the difference is
+7.2\% on Video Games and +2.7\% on Baby Products. The Baby margin over
SR-GNN is only .00138 and lies near observed seed variation, so we describe
it as a small numerical improvement rather than a broad dominance claim.

Diginetica produces the opposite ordering. \method{} exceeds six baselines, is
statistically tied with STAN, and ranks behind NARM. NARM is ahead by .01725,
corresponding to a relative deficit of 3.2\%. This bounds the claim: the
proposed architecture is strongest in the sparse, metadata-rich Amazon regime,
not universally.

\begin{table}[t]
\caption{Secondary ranking metrics for the strongest external comparator in
each domain.}
\label{tab:secondary}\centering\small
\begin{tabular}{lrrrrrr}
\toprule
Domain & Comparator & \method{} NDCG@20 & Comp. NDCG@20 & \method{} MRR@20 & Comp. MRR@20 \\
\midrule
Video Games & NARM & .06877 & .06316 & .04693 & .04264 \\
Baby Products & SR-GNN & .02492 & .02230 & .01699 & .01414 \\
Diginetica & NARM & .25589 & .26321 & .18159 & .18628 \\
\bottomrule
\end{tabular}
\end{table}

The same ordering appears on NDCG@20 and MRR@20 for the decisive Amazon
comparisons, while Diginetica again favours NARM. The conclusions therefore
do not depend on Recall@20 alone, even though Recall@20 remains the primary
selection target.

\subsection{Paired comparisons}
\begin{table}[t]
\caption{Headline paired comparisons against the strongest external baseline in
each domain. Intervals are query-clustered 95\% paired bootstrap ranges.}
\label{tab:headline-paired}\centering\small
\begin{tabular}{lrrr}
\toprule Domain & Strongest baseline & $\Delta$ Recall@20 & 95\% interval \\
\midrule
Video Games & NARM & +.00980 & [.00824,.01137] \\
Baby Products & SR-GNN & +.00138 & [.00038,.00236] \\
Diginetica & NARM & -.01725 & [-.02012,-.01435] \\
\bottomrule
\end{tabular}
\end{table}

The same paired machinery is used for favourable and unfavourable comparisons.
This matters because reporting only victories would hide the most important
boundary condition: on the domain with genuine anonymous sessions and weak
semantic content, a carefully tuned NARM remains superior.

Stratified analysis from the persisted rank artifacts asks where that ordering
changes. It does not. On Video Games, \method{} beats ID-only NARM on all
three target-popularity strata, with the largest Recall@20 gain on the tail
(+.01260 versus +.01057 on the head). On Baby Products, \method{} again beats
SR-GNN on head, torso, and tail targets, but only by .00099--.00291. The same
sign pattern holds across short, medium, and long context-length terciles on
both Amazon domains. Diginetica shows the opposite result on every slice:
NARM remains ahead on head, torso, and tail targets and on all context-length
terciles, with the largest deficit for \method{} on long contexts (-.01907).
Thus the Amazon advantage is broad within its regime, whereas the Diginetica
disadvantage is not concentrated in a single subset.

\subsection{What fusion contributes}
\begin{table}[t]
\caption{Constituent ablation under one validation-gated configuration and
three matched seeds.}
\label{tab:constituents}\centering\small
\begin{tabular}{lrrr}
\toprule Dataset & Memory only, $\beta=0$ & Neural only, $\beta=1$ & Validation-gated fusion \\
\midrule
Video Games & .11901 & .12571 & \textbf{.14685} \\
Baby Products & .03879 & .04873 & \textbf{.05346} \\
Diginetica & .49565 & .44852 & \textbf{.51681} \\
\bottomrule
\end{tabular}
\end{table}

Fusion improves over both endpoints on every domain. On Diginetica, the
selected system exceeds memory-only by .02116 with query-clustered 95\% CI
$[.01942,.02289]$, so the gain is smaller in interpretation than on Amazon but
not statistically ambiguous in the locked protocol. Unlike the
external-baseline comparison, this attribution is internally controlled: the
three columns share preprocessing, candidates, seeds, and selected component
configuration. The consistent result supports the central fusion claim even
where \method{} is not the strongest external system.

\subsection{Validation-gated broader expert fusion}
The constituent ablation above asks whether memory--neural fusion helps
relative to its own endpoints. A separate systems question is whether
\method{} should itself be fused with strong external memories. We therefore
evaluate a second rejection-complete family containing \method{}, STAN,
V-SKNN, NARM, and every RRF subset of size two or larger. Tmall omits NARM
because no matched checkpoint is available. The selector uses
$.5\,\mathrm{Recall}@10+.5\,\mathrm{Recall}@20$ on validation data; test
metrics are read only after the candidate is fixed.

\begin{table}[t]
\caption{Broader expert selection from persisted seed-42 rank artifacts.
$\Delta$R@20 is relative to the strongest individual expert. This diagnostic
is separate from the matched three-seed main comparison.}
\label{tab:expert-fusion}\centering\scriptsize
\setlength{\tabcolsep}{3.2pt}
\begin{tabular}{llrrrr}
\toprule
Dataset & Validation decision & R@10 & R@20 & NDCG@20 & $\Delta$R@20 \\
\midrule
Video Games & RRF(\method{}, STAN, NARM) & .10989 & .15354 & .07366 & +.00655 \\
Baby Products & RRF(\method{}, STAN, V-SKNN) & .04015 & .05872 & .02786 & +.00834 \\
Diginetica & RRF(STAN, NARM) & .41234 & .53980 & .26931 & +.00554 \\
Tmall & RRF(\method{}, STAN) & .23778 & .30751 & .16942 & +.01452 \\
\bottomrule
\end{tabular}
\end{table}

The decisions are more informative than a fourth average win. Video Games
admits STAN and NARM but rejects V-SKNN. Baby Products rejects NARM and keeps
the original three experts. Most importantly, Diginetica selects the control
RRF(STAN,NARM) and rejects \method{} itself, improving Recall@20 over
seed-42 NARM by 1.0\%. This converts a system-level loss into a successful
protocol decision without concealing which instantiation won. On
metadata-free Tmall, validation selects \method{}+STAN and improves Recall@20
over STAN by 5.0\%. Its NDCG@20
(.16942) remains below STAN (.17810), consistent with the recall-based
selection objective rather than a claim of metric-universal dominance.
Because these rows reuse one persisted rank artifact rather than three matched
seeds, they establish transfer behaviour and falsifiability, not significance.

\subsection{What metadata contributes}
The neural baselines use item identifiers, whereas the full Amazon PASGR
initialization uses item text. We therefore remove the semantic teacher and
rerun the complete gated system, then equalize semantic information for two
external baselines.

\begin{table}[t]
\caption{Amazon fairness rerun: external neural baselines with the same
semantic initialization used by PASGR. Values are three-seed means.}
\label{tab:semantic-fairness}\centering\small
\begin{tabular}{lrrr}
\toprule Dataset & GRU4Rec+sem & NARM+sem & \method{} \\
\midrule
Video Games & .13428 & \textbf{.15387} & .14685 \\
Baby Products & .05947 & \textbf{.06628} & .05346 \\
\bottomrule
\end{tabular}
\end{table}

\begin{table}[t]
\caption{Full system versus no-metadata rerun and strongest external baseline.}
\label{tab:nometa}\centering\small
\begin{tabular}{lrrr}
\toprule Dataset & Full \method{} & Without metadata & Strongest external baseline \\
\midrule
Video Games & .14685 & .13208 & NARM .13705 \\
Baby Products & .05346 & .05055 & SR-GNN .05208 \\
Diginetica & .51681 & .51756 (.00122) & NARM .53406 \\
\bottomrule
\end{tabular}
\end{table}

Without metadata, \method{} loses to the strongest Amazon baseline: -.00497
on Video Games with interval $[-.0069,-.0030]$, and -.00154 on Baby Products
with interval $[-.0027,-.0004]$. The fairness rerun goes further: against
semantic-initialized NARM, \method{} loses by -.00702 on Video Games
($[-.00847,-.00559]$) and -.01282 on Baby Products
($[-.01378,-.01185]$). The Amazon headline is therefore a system-level result
for a metadata-equipped method, not evidence that rank fusion alone defeats
every architecture under equal information.

Secondary metrics preserve the same fairness conclusion. On Video Games,
semantic-initialized NARM reaches NDCG@20/MRR@20 of .07112/.04809 against
\method{} at .06877/.04693; on Baby Products, the corresponding values are
.02914/.01891 against .02492/.01699. Metadata therefore changes not only the
hit-rate ordering but the broader top-rank quality picture.

On Diginetica, removing the weak semantic signal changes .51681 to
.51756$\pm$.00122, a reversal within seed noise. Both variants remain below
NARM. The Diginetica result is thus independent of useful side information,
while the Amazon advantage is not.

\subsection{Validation as a rejection mechanism}
\begin{table}[t]
\caption{Validation-selected neural switches.}
\label{tab:switches}\centering\small
\begin{tabular}{lrrrr}
\toprule Domain & Graph transport & Prototype transport & Contrastive loss & In-batch loss \\
\midrule
Video Games & .35 & Off & .15 & .10 \\
Baby Products & 0 & Off & .15 & 0 \\
Diginetica & .35 & Off & 0 & .10 \\
\bottomrule
\end{tabular}
\end{table}

Prototype transport is rejected on every domain. The agreement bonus is also
effectively null, changing Recall@20 by .00000 in separate ablation checks.
These negative results are not peripheral: they show that the selection
procedure can contradict the designer rather than merely accumulate modules.

Router families vary across seeds: Video Games and Diginetica alternate
between regime-level and bucketed routing, while Baby Products alternates
between bucketed and continuous routing. No domain selects one family in all
three seeds: the modal family appears in two of three runs for each domain,
and three of nine domain--seed decisions differ from their domain mode. We
therefore do not present fine-
grained routing as a standalone contribution. The defensible contribution is
the gate that may retain or discard routing granularity.

\subsection{Efficiency}
\begin{table}[t]
\caption{Representative cost profile.}
\label{tab:efficiency}\centering\small
\begin{tabular}{lr}
\toprule Operation & Measured cost \\
\midrule
Memory index construction & 1.2--1.8 s, single-process CPU \\
Three memories plus fusion & 3.1--3.6 ms/query, CPU \\
PASGR inference & about .3 ms CPU; .05 ms Metal backend \\
Selected PASGR training & about 5 min \\
Selected \method{} reproduction & about 6 min \\
One-off 16-cell neural search & about 80 min per domain and seed \\
One baseline training run & 5--15 min \\
\bottomrule
\end{tabular}
\end{table}

The selected system is comparable to one neural baseline run and requires no
GPU at inference. The one-off search is substantially larger and is reported
separately. These costs position \method{} between fully training-free
retrieval and learned ensembles.

\section{Discussion}
The results support a simple diagnosis. Sparse recommendation does not lack
candidate signals; it lacks a reliable rule for deciding which signal deserves
influence. Memory dominates when local repetition and transition structure are
sufficient. Neural evidence helps when semantic or sequential generalization
adds information not already present in those memories. Rank-space fusion
removes the irrelevant question of how to calibrate incomparable raw scores,
and validation gating addresses the substantive question of whether the neural
residual helps in the target regime.

This distinction also locates the novelty precisely. Weighted reciprocal-rank
fusion, neighbourhood retrieval, and recurrent encoders are established
ingredients. \method{} contributes neither another score formula nor a claim
that these ingredients dominate NARM. It contributes a rejection-complete
decision protocol and an empirical diagnosis produced by that protocol. The
strong NARM results are therefore informative counterexamples: they delimit
the regime in which a memory-first residual is justified rather than
invalidating the method's central claim.

The protocol also reports when the instantiated system loses, rather than
retaining only configurations or comparisons that produce a win. The
semantic-matched NARM reversal on Amazon and the ID-only NARM advantage on
Diginetica are therefore outputs of the same evaluation contract as the
positive endpoint-fusion results.

This interpretation aligns with current evidence from two directions. First,
shortcut audits show that expensive sequential modeling is often credited for
structure already recoverable by simple graphs~\cite{graphshortcut2026}.
Second, model-selection research shows meaningful complementarity but weak
per-instance selection~\cite{algselect2025}. \method{} responds by making
memory explicit, keeping routing optional, and requiring held-out evidence
before complexity is admitted.

\section{Limitations}
The first limitation is comparison symmetry. Rich semantic initialization is
available to \method{} and to the fairness-rerun baselines, but not to every
external model in the main table. Accordingly, the Amazon result is stated as
the best result among the evaluated ID-only external baselines, not proof that
fusion is superior under every information budget.

Second, the Baby Products margin over SR-GNN is small. Third, NARM
significantly outperforms \method{} on Diginetica, showing that a strong
neural primary model remains preferable in some dense session regimes. Fourth,
router-family variation across seeds prevents a universal claim about
per-query routing. Fifth, the multi-expert extension is not monotone across
domains: it excludes V-SKNN on Video Games, excludes NARM on Baby Products,
and excludes \method{} on Diginetica. Sixth, the
approximately seventy dependent
validation decisions create selection uncertainty even though all grids are
fixed and the test split is isolated. A simple Bonferroni-style view would
treat this search as too wide for isolated marginal validation differences;
consequently, only stable cross-domain or paired-test effects are used for
central claims. Seventh, broader fusion with external memories is currently
based on one persisted rank artifact per domain and lacks the matched-seed
protocol of the main comparison; Tmall also lacks a completed matched-seed
neural-baseline suite.
Finally, the RRF(STAN,NARM) control is currently a single-artifact diagnostic;
matched-seed replication is necessary before making a significance claim about
the selected expert family.

\section{Conclusion}
\method{} combines three explicit memories with a compact neural residual
entirely in rank space. The design removes score-scale mismatch, allows
$\beta=0$ as a training-free boundary, and assigns every discrete decision to
validation rather than to test inspection. Across two sparse Amazon domains
and Diginetica, fusion improves over both of its own constituents. The
complete metadata-equipped system ranks first on both Amazon domains, while
NARM remains significantly stronger on Diginetica.

The most transferable result is methodological. Validation is useful not only
for choosing a weight, but for rejecting unnecessary complexity: prototype
transport is removed, the agreement bonus is inert, and unstable fine-grained
routing is kept outside the central claim. The broader expert diagnostic
provides a second rejection test: Video Games excludes V-SKNN, Baby Products
excludes NARM, and Diginetica excludes \method{} itself in favour of
RRF(STAN,NARM). At the same time, the no-metadata and semantic-fairness
analyses prevent semantic information from being mistaken for fusion quality.

The broader lesson is more conservative than a single global winner: sparse
recommendation benefits from validation-gated adaptation across information
regimes, but the right endpoint may be a memory-first fusion, a stronger
neural primary model, or a simpler ensemble of external experts.

\paragraph{AI assistance disclosure.}
OpenAI Codex was used to assist with code review, experiment orchestration,
statistical-procedure checks, and language editing. The authors remain fully
responsible for the experiments, claims, accuracy, integrity, and originality
of the manuscript.

\bibliographystyle{splncs04}
\bibliography{references}
\end{document}
